% GENERATED FILE. Edit canonical lessons, then run npm run generate:books.
% canonical-source-hash: fnv1a64:b82f5c20ee1d32b4
% canonical-lessons: PT-C453-talao, PT-C453-pechincha, PT-C453-liquidacao, PT-C453-regatear, PT-C453-embrulhar

\chapter{Na caixa}
\label{ch:pt-na-caixa}
\hlchaptermodality{453}

\begin{chapteropening}
\textbf{By the end of this chapter:} \emph{I can pay and wrap things: the till receipt, a bargain, a clearance sale, haggling and gift-wrapping.}

Complete the last lesson of chapter 453: I can pay and wrap things: the till receipt, a bargain, a clearance sale, haggling and gift-wrapping.
\end{chapteropening}

\section[o talão]{o talão — the till receipt, the slip; the stub}

\label{lesson:PT-C453-talao}



\begin{quote}
Before the new one: say the Portuguese for the optician, then the Portuguese for the stain.
\end{quote}

\subsection*{You'll want to know: o talão}
\textbf{o talão} — \textquotedblleft{}the till receipt, the slip; the stub\textquotedblright{}.

\begin{quote}
Guardei o talão para trocar a camisa.
\end{quote}

\begin{quote}
\emph{I kept the receipt so I could exchange the shirt.}
\end{quote}

\emph{O talão} is the slip from the till; \emph{a fatura} has your tax number.

\begin{cousinweb}
\emph{talão} was borrowed from French \emph{talon} 'heel; stub of a booklet' (EtymDB), from Vulgar Latin \emph{tālō}, on Latin \emph{tālus} 'ankle, heel'.
\end{cousinweb}

\begin{grammarlens}[title={What you've built}]
One more word for telling, explaining and describing at B1.
\end{grammarlens}

\subsection*{Your turn}
\begin{itemize}
\raggedright
  \item \emph{Say it:} \emph{o talão}
  \item \emph{Say it:} \emph{o talão}, once more
  \item \emph{Say it:} a sentence of your own with \emph{o talão}
  \item \emph{Recall:} say the Portuguese for the optician, then the Portuguese for the stain, then say \emph{o talão} again
\end{itemize}

\begin{tcolorbox}[breakable,colback=teal!4,colframe=teal!35!black,title={Before you move on}]
What does o talão mean? (\textquotedblleft{}The till receipt, the slip; the stub\textquotedblright{}.) Say it once more.
\end{tcolorbox}
\section[a pechincha]{a pechincha — the bargain, the snip}

\label{lesson:PT-C453-pechincha}



\begin{quote}
Before the new one: say the Portuguese for the stain, then the Portuguese for the till receipt.
\end{quote}

\subsection*{You'll want to know: a pechincha}
\textbf{a pechincha} — \textquotedblleft{}the bargain, the snip\textquotedblright{}.

\begin{quote}
Este casaco foi uma pechincha.
\end{quote}

\begin{quote}
\emph{This coat was a bargain.}
\end{quote}

\emph{Que pechincha!} — what a bargain!

\begin{cousinweb}
The origin of \emph{pechincha} is not settled, and EtymDB has no entry for it.
\end{cousinweb}

\begin{grammarlens}[title={What you've built}]
One more word for telling, explaining and describing at B1.
\end{grammarlens}

\subsection*{Your turn}
\begin{itemize}
\raggedright
  \item \emph{Say it:} \emph{a pechincha}
  \item \emph{Say it:} \emph{a pechincha}, once more
  \item \emph{Say it:} a sentence of your own with \emph{a pechincha}
  \item \emph{Recall:} say the Portuguese for the stain, then the Portuguese for the till receipt, then say \emph{a pechincha} again
\end{itemize}

\begin{tcolorbox}[breakable,colback=teal!4,colframe=teal!35!black,title={Before you move on}]
What does a pechincha mean? (\textquotedblleft{}The bargain, the snip\textquotedblright{}.) Say it once more.
\end{tcolorbox}
\section[a liquidação]{a liquidação — the clearance sale; the settling (of a debt)}

\label{lesson:PT-C453-liquidacao}



\begin{quote}
Before the new one: say the Portuguese for the till receipt, then the Portuguese for the bargain.
\end{quote}

\subsection*{You'll want to know: a liquidação}
\textbf{a liquidação} — \textquotedblleft{}the clearance sale; the settling (of a debt)\textquotedblright{}.

\begin{quote}
A loja está em liquidação total.
\end{quote}

\begin{quote}
\emph{The shop is having a closing-down sale.}
\end{quote}

\emph{Os saldos} are the seasonal sales.

\begin{cousinweb}
\emph{liquidação} is built on \emph{liquidar} 'to settle, to clear', a learned formation on Latin \emph{liquidus} 'liquid, clear' (Lewis \& Short): a debt made clear.
\end{cousinweb}

\begin{grammarlens}[title={What you've built}]
One more word for telling, explaining and describing at B1.
\end{grammarlens}

\subsection*{Your turn}
\begin{itemize}
\raggedright
  \item \emph{Say it:} \emph{a liquidação}
  \item \emph{Say it:} \emph{a liquidação}, once more
  \item \emph{Say it:} a sentence of your own with \emph{a liquidação}
  \item \emph{Recall:} say the Portuguese for the till receipt, then the Portuguese for the bargain, then say \emph{a liquidação} again
\end{itemize}

\begin{tcolorbox}[breakable,colback=teal!4,colframe=teal!35!black,title={Before you move on}]
What does a liquidação mean? (\textquotedblleft{}The clearance sale; the settling (of a debt)\textquotedblright{}.) Say it once more.
\end{tcolorbox}
\section[regatear]{regatear — to haggle}

\label{lesson:PT-C453-regatear}



\begin{quote}
Before the new one: say the Portuguese for the bargain, then the Portuguese for the clearance sale.
\end{quote}

\subsection*{You'll want to know: regatear}
\textbf{regatear} — \textquotedblleft{}to haggle\textquotedblright{}.

\begin{quote}
Na feira, as pessoas regateiam o preço.
\end{quote}

\begin{quote}
\emph{At the market, people haggle over the price.}
\end{quote}

In shops with fixed prices nobody \emph{regateia}.

\begin{cousinweb}
The origin of \emph{regatear} is not settled; it is linked with \emph{regatão} 'street trader', and EtymDB has no entry for either word.
\end{cousinweb}

\begin{grammarlens}[title={What you've built}]
One more word for telling, explaining and describing at B1.
\end{grammarlens}

\subsection*{Your turn}
\begin{itemize}
\raggedright
  \item \emph{Say it:} \emph{regatear}
  \item \emph{Say it:} \emph{regatear}, once more
  \item \emph{Say it:} a sentence of your own with \emph{regatear}
  \item \emph{Recall:} say the Portuguese for the bargain, then the Portuguese for the clearance sale, then say \emph{regatear} again
\end{itemize}

\begin{tcolorbox}[breakable,colback=teal!4,colframe=teal!35!black,title={Before you move on}]
What does regatear mean? (\textquotedblleft{}To haggle\textquotedblright{}.) Say it once more.
\end{tcolorbox}
\section[embrulhar]{embrulhar — to wrap (up)}

\label{lesson:PT-C453-embrulhar}



\begin{quote}
Before the new one: say the Portuguese for the clearance sale, then the Portuguese for to haggle.
\end{quote}

\subsection*{You'll want to know: embrulhar}
\textbf{embrulhar} — \textquotedblleft{}to wrap (up)\textquotedblright{}.

\begin{quote}
A senhora embrulhou o presente em papel azul.
\end{quote}

\begin{quote}
\emph{The woman wrapped the present in blue paper.}
\end{quote}

\emph{É para oferta?} — Is it a gift? Then they \emph{embrulham} it. \emph{Um embrulho} — a parcel.

\begin{cousinweb}
\emph{embrulhar} comes from Old Portuguese \emph{envurullar} (EtymDB); its earlier origin is not settled.
\end{cousinweb}

\begin{grammarlens}[title={What you've built}]
One more word for telling, explaining and describing at B1.
\end{grammarlens}

\subsection*{Your turn}
\begin{itemize}
\raggedright
  \item \emph{Say it:} \emph{embrulhar}
  \item \emph{Say it:} \emph{embrulhar}, once more
  \item \emph{Say it:} a sentence of your own with \emph{embrulhar}
  \item \emph{Recall:} say the Portuguese for the clearance sale, then the Portuguese for to haggle, then say \emph{embrulhar} again
\end{itemize}

\begin{tcolorbox}[breakable,colback=teal!4,colframe=teal!35!black,title={Before you move on}]
What does embrulhar mean? (\textquotedblleft{}To wrap (up)\textquotedblright{}.) Say it once more.
\end{tcolorbox}
